% Generated by scripts/generate-worked-calculations.R; do not edit by hand.
% Registry ID: reg.annual.expectation_interval
% Source notebook: notebooks/support/regression-analysis/support.Rmd
For the \emph{Case: US SteamCo}, the annual expectation is obtained by summing the 12 monthly expectations after adjusting the March 2014 expectation for the corroborated winter storm effect.\footnote{See Exercise \ref{ex:Combining_12_monthly_predictions}} The resulting annual expectation is:

\begin{equation*}
\hat{y}_{\text{total}} = 190,680,684.
\end{equation*}

Using the variances and covariance of the estimated regression coefficients together with the residual variance yields a combined prediction standard error of

\begin{equation*}
s_{\text{comb}} = 6,323,301.
\end{equation*}

The corresponding annual prediction-variance components are approximately $14,886,605,275,947.4648$ from coefficient uncertainty and $25,097,531,113,309.0156$ from future residual variation. Thus, about 37.2\% of annual prediction variance comes from parameter estimation and 62.8\% from future-period noise.

With a 99\% prediction interval and the appropriate $t$ critical value, the resulting acceptable prediction range for the annual revenue is

\begin{equation*}
173,291,634
\le
y_{\text{total}}
\le
208,069,735.
\end{equation*}

The recorded annual revenue for 2014 equals

\begin{equation*}
y_{\text{recorded}} = 174,778,400,
\end{equation*}
so that the annual difference, defined as recorded revenue minus expected revenue, is

\begin{equation*}
174,778,400 - 190,680,684 = -15,902,284.
\end{equation*}

The annual difference is consistent with the uncertainty measured by the adjusted 99\% acceptable difference range. The substantive analytical procedure therefore does not identify a remaining significant unexplained difference at the annual level, and the auditor can proceed to determine whether the evidence and achieved assurance are sufficient to reject material misstatement.

Consequently, the annual difference of 15,902,284 is first compared with the adjusted 99\% acceptable difference range. Because the annual difference falls within that range, the substantive analytical procedure does not identify a significant unexplained difference requiring further investigation.
